El problema de llevar \emph{software} de una computadora a otra y que siga funcionando (proceso conocido como Despliegue de \emph{Software}) es un yugo al que las Ciencias de la Computación no han dado una solución definitiva, esencialmente porque para que una pieza de \emph{software} funcione correctamente no solamente depende del código fuente en el que está escrita, sino del contexto que lo rodea (\emph{hardware}, sistema operativo, dependencias, etc.). Bajo ese contexto surge Nix como un manejador de paquetes y sistema de construcción\footnote{Sistemas que automatizan la ejecución de tareas repetitivas, usualmente para crear artefactos de \emph{software} que pueden ser desplegados} que utiliza un enfoque inspirado en la pureza funcional, donde cada paquete define explícitamente el contexto en el que espera ser construido y ejecutado. Este enfoque ha demostrado poder crear más de 700 mil artefactos binariamente idénticos en diferentes computadoras y contar con uno de los repositorios de paquetes más grandes a la fecha de este documento. La idea evolucionó al punto de definir el estado casi completo de un sistema operativo mediante un solo archivo de configuración.

Sin embargo, a pesar de sus capacidades prometedoras, Nix no ha gozado de la misma adopción que otras herramientas que abordan los mismos problemas de reproducibilidad como Docker, Conda o VirtualEnv u otros manejadores de paquetes. Las causas de acuerdo a la comunidad son varias: documentación compleja, errores crípticos, o un lenguaje de programación difícil de dominar; muchos de ellos no siendo problemas técnicos sino de experiencia de uso, y que caen en el rango de estudio del ``DX'' (Experiencia de Desarrollo por sus siglas en inglés). Este trabajo se enfoca en verificar si proveer una nueva forma de interactuar con Nix a través de un lenguaje de propósito general como TypeScript puede reducir la barrera de entrada para nuevos usuarios y por ende mejorar su DX.
